% [原文 p.310]
BPS 态，它们按 $\operatorname{Spin}(8)$ 的八维表示变换，而该表示由 $(n_e,n_m)$ 模 $2$ 约化确定。

$W^\pm$-玻色子的低能电荷为 $(2,0)$。它们衰变为基本超多重态的过程是边缘稳定的。S-对偶要求存在电荷为 $(2n_e,2n_m)$ 的 BPS 态，其中 $n_e,n_m$ 取遍所有互素整数。

对规范群为 $SU(2)$、$F=4$ 的无质量 $4d$ $\mathcal{N}=2$ SYM，其 Seiberg–Witten 曲线具有三对偶性与 $\SL_2(\Z)$-不变性，这为在该理论中实现 S-对偶提供了有力证据。

当基本超多重态有质量时，标度不变性与 $U(1)_R$ 对称性都被明显破坏。超多重态的质量按味对称群的伴随表示变换，因此三对偶群 $\mathfrak{S}_3$ 如下作用在它们上面 [314, 第 10 节]。交换旋量表示 $s$ 与 $c$ 而保持 $v$ 不变的那个变换，其作用是翻转 $m_4$ 的符号，即 $m_4\to-m_4$，而保持其余质量 $m_1,m_2,m_3$ 不变。交换矢量 $v$ 与旋量 $s$ 而保持 $c$ 不变的那个变换，作用为
\begin{equation}
\begin{aligned}
m_1&\mapsto \tfrac{1}{2}(m_1+m_2+m_3+m_4)\,,\\
m_2&\mapsto \tfrac{1}{2}(m_1+m_2-m_3-m_4)\,,\\
m_3&\mapsto \tfrac{1}{2}(m_1-m_2+m_3-m_4)\,,\\
m_4&\mapsto \tfrac{1}{2}(m_1-m_2-m_3+m_4)\,.
\end{aligned}\tag{20.14}
\end{equation}
这两个变换一起生成整个三对偶群 $\mathfrak{S}_3$。

% [Tutor Note: 式 (20.14) 可核对如下：该 $4\times4$ 矩阵是 $\frac12$ 倍的 Hadamard 型矩阵，各行两两正交且范数为 $1$，故它是正交变换（与「交换 $v$ 与 $s$、保持 $c$」的三对偶自同构相符）；它自身满足 $T^2=1$，而 $m_4\mapsto-m_4$ 记为 $F$ 也满足 $F^2=1$。验证乘积：$TF$ 的矩阵为 $\frac12\begin{pmatrix}1&1&1&-1\\1&1&-1&1\\1&-1&1&1\\-1&1&1&1\end{pmatrix}$，直接计算得 $(TF)^3=\mathrm{Id}$，故 $\langle T,F\rangle$ 是 6 阶群，即 $\mathfrak{S}_3$——与原文「这两个变换生成整个三对偶群」一致。]
\subsection{穿墙}

对带 $1\le F\le 3$ 个无质量基本味的 $4d$ $\mathcal{N}=2$ $SU(2)$ SYM，其 Seiberg–Witten 解的细致分析表明，这些理论中会发生穿墙 [58]。边缘稳定性的墙仍然同胚于一个圆周，并且穿过模空间 $\mathcal{C}^{(F)}$ 强耦合区域中的全部奇点：

\begin{itemize}
\item 对 $F=1$，在 $\mathcal{C}^{(1)}$ 中墙的外侧，BPS 谱由电荷为 $\pm(n_e,1)$（对所有 $n_e\in\Z$）的双子、基本超多重态 $\pm(1,0)$ 以及 $W^\pm$-玻色子 $\pm(2,0)$ 组成。边缘稳定性曲线穿过强耦合区域中的三个奇点。在墙的内侧，BPS 谱约化为造成强耦合奇点的那些态：磁单极子 $\pm(0,1)$ 以及双子 $\pm(-1,1)$ 与 $\pm(-2,1)$。
\item 当 $F=2$ 时，弱耦合 BPS 谱由所有双子 $\pm(n_e,1)$（对所有 $n_e\in\Z$）、基本超多重态 $\pm(1,0)$ 与 $W^\pm$-玻色子 $\pm(2,0)$ 组成。

% [原文 p.311]
在墙的内侧，BPS 谱只由磁单极子 $\pm(0,1)$ 与双子 $\pm(1,1)$ 组成。
\item 当 $F=3$ 时，弱耦合 BPS 谱包括所有双子 $\pm(n_e,1)$ 与 $\pm(2n+1,2)$（对所有 $n_e,n\in\Z$），连同基本超多重态 $\pm(1,0)$ 与 $W^\pm$-玻色子 $\pm(2,0)$。在墙的内侧，BPS 谱只由磁单极子 $\pm(0,1)$ 与双子 $\pm(-1,2)$ 组成。
\end{itemize}

\begin{remark}[20.27]
在这三种情形中，穿墙都调和了 S-对偶与如下事实：BPS 谱在由单值变换所生成的群下并不保持不变。当穿过一面墙时——典型地是在绕强耦合奇点时——某些 BPS 态会衰变。
\end{remark}

当基本超多重态有质量时，奇点与边缘稳定性墙都在模空间 $\mathcal{C}^{(F)}$ 内发生位移。例如，若第一个超多重态的裸质量 $m_1\gg\Lambda_F$，则在弱耦合区域中存在一个位于 $u\simeq m_1^2/2$ 的奇点，该基本超多重态在那里变为无质量。相应的单值变换由式 (20.12) 给出，或等价地：
\begin{equation}
\begin{pmatrix}n_e\\ n_m\\ S\end{pmatrix}\mapsto
\begin{pmatrix}1&-1&0\\0&1&0\\0&1&1\end{pmatrix}
\begin{pmatrix}n_e\\ n_m\\ S\end{pmatrix}
=\begin{pmatrix}n_e-n_m\\ n_m\\ S+n_m\end{pmatrix}.\tag{20.15}
\end{equation}
在弱耦合区域，BPS 态的味荷 $S$ 是有界的。把式 (20.15) 应用于电荷为 $(0,1,1)$ 的磁单极子可以看出，存在一面穿过该奇点的边缘稳定性墙。因此，在基本超多重态质量很大的 $4d$ $\mathcal{N}=2$ $SU(2)$ SYM 理论中，即使在库仑分支的弱耦合区域也会发生穿墙。

\section{Seiberg–Witten 分析：一般情形}

Seiberg 与 Witten 的分析已被迅速推广到规范代数为各种情形的 $4d$ $\mathcal{N}=2$ 纯 SYM 理论，包括 $\mathfrak{su}(N)$ [25, 232]、$\mathfrak{so}(2N+1)$ [100]、$\mathfrak{so}(2N)$ [64]，以及带基本物质的 $\mathcal{N}=2$ SQCD 理论 [185, 28, 30, 184]。

如同 20.1 与 20.2 节中详细讨论的 $\mathfrak{su}(2)$ 理论，这些理论的精确非微扰动力学可以完全由一族称为 Seiberg–Witten 曲线的黎曼曲面来刻画。这些曲线由理论的库仑分支参数化，并带有称为 Seiberg–Witten 微分的亚纯微分，它定义到相差恰当 1-形式。其背后的数学结构是代数完全可积系统 [114]，我们稍后讨论。

在一般的渐近自由或共形 $4d$ $\mathcal{N}=2$ 量子场论中，库仑分支 $\mathcal{C}$ 是一个复 $r$ 维仿射空间，其中 $r$ 是理论的秩。

% [原文 p.312]
当理论是规范代数为 $\mathfrak{g}$ 的拉格朗日 SYM 理论时，$\mathcal{C}$ 上的自然坐标来自伴随标量 $\phi$ 的 $\mathfrak{g}$-不变多项式，如 18.2 节所述。例如，当 $\mathfrak{g}=\mathfrak{su}(N)$ 时秩为 $r=N-1$，库仑分支 $\mathcal{C}$ 可以由 $u_2=\langle\tr\phi^2\rangle$、$u_3=\langle\tr\phi^3\rangle$、……、$u_N=\langle\tr\phi^N\rangle$ 参数化。

在 $\mathcal{C}$ 的一般点处，低能理论是 $4d$ $\mathcal{N}=2$ 广义 Maxwell 理论，规范群为 $\U(1)^r$。理论还可能有阿贝尔味对称性 $\U(1)^f$。低能电荷为
\[
(n_e^1,n_{m,1},\dots,n_e^r,n_{m,r},S_1,\dots,S_f),
\]
它们构成一个格 $\Gamma\simeq\Z^{2r+f}$，带有 Dirac–Schwinger–Zwanziger（DSZ）配对：
\[
\Big\langle (n_e^i,n_{m,i},S_k)\,,\big((n_e^i)',n'_{m,i},S'_k\big)\Big\rangle
=\Big(\sum_{i=1}^{r}\Big)n_e^i n'_{m,i}-n_{m,i}(n_e^i)'.
\]

对每个 $i=1,\dots,r$，第 $i$ 个 $\mathcal{N}=2$ 阿贝尔矢量超场的真空期望值（VEV）$a_i$ 是 $\mathcal{C}$ 上的局部全纯函数。一般地，集合 $(a_1,\dots,a_r)$ 构成 $\mathcal{C}$ 上的一组局部全纯坐标。$\mathcal{C}$ 上的特殊凯勒度量可以用低能阿贝尔规范耦合 $\tau^{ij}$ 表示为
\begin{equation}
\dd s^2=\operatorname{Im}\tau^{ij}\dd a_i\dd\bar a_j.\tag{20.16}
\end{equation}
矩阵 $(\tau^{ij})$ 由于 $\mathcal{N}=2$ 超对称而是对称的。这又意味着存在局部全纯的预势 $\mathcal{F}$，使得
\[
\tau^{ij}=\frac{\partial^2\mathcal{F}}{\partial a_i\partial a_j}.
\]
此外，定义
\[
a_D^i=\frac{\partial\mathcal{F}}{\partial a_i},
\]
其中 $a_i$ 与 $a_D^i$ 称为特殊凯勒流形 $\mathcal{C}$ 上的特殊坐标。有
\[
\dd s^2=\operatorname{Im}\dd a_D^i\dd\bar a_i.
\]
$4d$ $\mathcal{N}=2$ 中心荷为
\begin{equation}
\mathcal{Z}=\Big(\sum_{i=1}^{r}n_e^i a_i+n_{m,i}a_D^i\Big)+\Big(\sum_{k=1}^{f}S_k\frac{m_k}{\sqrt{2}}\Big),\tag{20.17}
\end{equation}
其中 $m_k$ 是与味荷相伴的常数。任一单粒子态的质量 $M$ 都受 BPS 不等式约束：
\[
M\ge|\mathcal{Z}|,
\]

% [原文 p.313]
等号成立当且仅当该态是 BPS 态。给定单粒子态的电荷对应于电荷格中一个点 $\gamma\in\Gamma$，但并非所有 $\gamma$ 都必然对应 BPS 态。

\begin{definition}[20.28]
理论的 BPS 谱是使 BPS 态存在的那些电荷 $\gamma$ 构成的集合，连同这些态的重数。
\end{definition}

我们将在下一章讨论的谱网络，是确定许多 $4d$ $\mathcal{N}=2$ 理论 BPS 谱的一个非常有力的方法。

在 $\mathcal{C}$ 上度量奇点的附近，特殊坐标与常数 $m_k$ 会经历单值变换：
\begin{equation}
(a_1\ a_D^1\ \cdots\ a_r\ a_D^r\ m_1/\sqrt{2}\ \cdots\ m_f/\sqrt{2})
\begin{pmatrix} M & 0_{2r\times f}\\ C & \mathrm{Id}_f\end{pmatrix},\tag{20.18}
\end{equation}
其中 $M\in\Sp(2r,\Z)$，以保证 DSZ 配对被保持，而 $C\in\mathcal{M}_{f\times 2r}(\Z)$。电、磁与味荷按这类单值矩阵的逆作变换：
\[
\begin{pmatrix}
n_e^1\\ n_{m,1}\\ \vdots\\ n_r^1\\ n_{m,r}\\ S_1\\ \vdots\\ S_f
\end{pmatrix}
\mapsto
\begin{pmatrix} M^{-1} & 0_{2r\times f}\\ \widetilde{C} & \mathrm{Id}_{k\times f}\end{pmatrix}
\begin{pmatrix}
n_e^1\\ n_{m,1}\\ \vdots\\ n_r^1\\ n_{m,r}\\ S_1\\ \vdots\\ S_f
\end{pmatrix},
\]
从而保持中心荷 $\mathcal{Z}$。这一形式表明：味荷 $S_k$ 可以获得来自阿贝尔规范荷的贡献，但在这些变换下规范荷不能被味荷改变。

% [Tutor Note: 两处核对/存疑。（i）DSZ 配对式中 $\big(\sum_{i=1}^r\big)$ 的括号在原文里只包住了求和号，应是排版疏漏；按配对的定义应读作 $\sum_{i=1}^{r}\big(n_e^i n'_{m,i}-n_{m,i}(n_e^i)'\big)$。核对方法：把两个电荷互换得 $\sum_i(n'^i_e n_{m,i}-n'_{m,i}n_e^i)$，恰为原式的相反数，故配对反对称，与 $M\in\Sp(2r,\Z)$ 保持配对的断言相容。（ii）电荷列向量第四行原文印作 $n_r^1$，按上下文应为 $n_e^r$（与首行 $n_e^1$ 及式 (20.18) 中 $a_1,a_D^1,\dots,a_r,a_D^r$ 的顺序对照可见）；同理分块矩阵中的 $\mathrm{Id}_{k\times f}$ 应为 $\mathrm{Id}_f$（对照式 (20.18) 的 $\mathrm{Id}_f$，味荷分块是 $f\times f$）。此处均照原文译出，读者按 $n_e^r$、$\mathrm{Id}_f$ 理解即可。]
\subsection{Seiberg–Witten 解}

在规范代数具有任意秩的各种 $4d$ $\mathcal{N}=2$ SYM 理论中——无论是纯理论 [25, 232, 100, 64] 还是带基本物质的理论 [185, 28, 30, 184]——已经证明，库仑分支 $\mathcal{C}$ 上的特殊坐标 $a_i$ 与 $a_D^i$ 可以表为某个亚纯微分 $\lambda$ 在某些超椭圆曲线上的周期积分。这把原始的 $\mathfrak{su}(2)$ 规范理论的 Seiberg–Witten 分析推广到更高秩的规范代数，从而给出这些理论中量子动力学的有力而精确的非微扰描述。

% [原文 p.314]
记 $\mathbf{u}=(u_1,\dots,u_r)$ 为库仑分支算子，即 $\mathcal{C}$ 上的规范不变坐标，则有：

\begin{proposition}[20.29]
每个这样的 $4d$ $\mathcal{N}=2$ 理论都对应于由 $\mathcal{C}$ 参数化的一族全纯黎曼曲面 $(\Sigma_{\mathbf{u}})$，称为 Seiberg–Witten 曲线。这些曲面由一个多项式方程定义，该方程涉及复变量 $z\in\C\PP^1$、$x\in\C$、库仑分支坐标 $\mathbf{u}$ 以及超多重态的裸质量 $\mathbf{m}$，后者被视为理论的固定参数。
\end{proposition}

在库仑分支上固定的一个一般点处，即对 $\mathbf{u}$ 的一个固定一般值，$\C\PP^1_z\times\C_x$ 中的黎曼曲面 $\Sigma_{\mathbf{u}}$ 是亏格为 $r$、带 $f$ 个穿孔的曲面。此外，$\Sigma_{\mathbf{u}}$ 是 $\C\PP^1_z$ 的 $r$ 次分歧覆叠。当 $\mathbf{u}$ 变化时，曲面可在库仑分支的非一般点 $\mathbf{u}_0$ 处退化，这解释了 $\mathcal{C}$ 上的度量奇点。最简单的退化类型发生在覆叠 $\Sigma_{\mathbf{u}}\to\C\PP^1_z$ 的两个分支点相撞时。这类退化称为节点奇点。

\begin{proposition}[20.30]
特殊坐标 $a_i$、$a_D^i$ 与质量 $m_k$ 由亚纯微分 $\lambda_{\mathbf{u}}$（称为 Seiberg–Witten 微分）在把穿孔紧化后的 $\Sigma_{\mathbf{u}}$ 上的周期积分给出。$\lambda_{\mathbf{u}}$ 的极点位于 $\Sigma_{\mathbf{u}}$ 的穿孔处。在 $\mathbf{u}$ 的固定值处，同调群 $H_1(\Sigma_{\mathbf{u}},\Z)$ 构成秩为 $2r+f$ 的格。在 $\C$ 上，该同调群构成格的局部系统，带有反对称相交配对 $(\cdot,\cdot)$，其核中包含绕穿孔的圈。
\end{proposition}

\begin{definition}[20.31]
$H_1(\Sigma_{\mathbf{u}},\Z)$ 的一个辛基由 $\Sigma_{\mathbf{u}}$ 上的圈 $A_1,B_1,\dots,A_r,B_r,C_1,\dots,C_f$ 组成，其中圈 $C_k$ 围绕穿孔，并且相交配对满足
\[
(A_i,B_j)=-(B_j,A_i)=\delta_{ij},
\]
其余相交全为零。
\end{definition}

在 $\mathcal{C}$ 上局部地，可以选取 $H_1(\Sigma_{\mathbf{u}},\Z)$ 的一个辛基，从而把带有相交配对 $(\cdot,\cdot)$ 的同调格局部系统 $H_1(\Sigma_{\mathbf{u}},\Z)$ 与带有 DSZ 配对 $\langle\cdot,\cdot\rangle$ 的电荷格局部系统 $\Gamma_{\mathbf{u}}$ 等同起来。于是：

\begin{proposition}[20.32]
对每个 $i=1,\dots,r$，特殊坐标 $a_i$ 与 $a_D^i$ 由 Seiberg–Witten 微分 $\lambda$ 沿相应同调圈的周期积分给出：
\[
a_i=\oint_{A_i}\lambda,\qquad\text{以及}\qquad a_D^i=\oint_{B_i}\lambda .
\]
味参数 $m_k/\sqrt{2}$ 由 Seiberg–Witten 微分 $\lambda$ 在 $\Sigma_{\mathbf{u}}$ 的穿孔处的留数得到。
\end{proposition}

% [原文 p.315]
纤维化 $\Sigma\to\mathcal{C}$ 中的奇点在同调群 $H_1(\Sigma_{\mathbf{u}},\Z)$ 中诱导单值变换，其形式如前所述，并可由 Picard–Lefschetz 单值公式得到。这些单值变换给出式 (20.18) 中描述的电荷格的单值变换。

一个重要的约束来自如下事实：规范耦合 $\tau^{ij}=\partial a_D^i/\partial a_j$ 必须非退化，且其虚部必须正定，以保证式 (20.16) 中库仑分支上的度量有良好定义且正定。黎曼曲面 $\Sigma_{\mathbf{u}}$ 的紧化是紧的且亏格为 $r$，因此它们容许 $r$ 个独立的全纯（阿贝尔）微分 $\omega_1,\dots,\omega_r$。对 $H_1(\Sigma_{\mathbf{u}})$ 的一个辛基 $A_1,B_1,\dots,A_r,B_r$ 以及周期矩阵
\[
u_{ij}=\oint_{A_i}\omega_j,\qquad\text{以及}\qquad v_{ij}=\oint_{B_i}\omega_j,
\]
Riemann 双线性关系蕴含矩阵 $(v_{ij}/u_{ij})$ 非退化且虚部正定。因此，我们要求导数
\[
\frac{\partial\lambda}{\partial u_i}
\]
构成 $\Sigma_{\mathbf{u}}$ 上阿贝尔微分的一组基。

\begin{remark}[20.33]
Seiberg–Witten 微分 $\lambda$ 不是唯一确定的：它仍可相差一个恰当微分的平移，即 $\lambda\mapsto\lambda+\dd\alpha$，因为这不会影响 $\lambda$ 的周期。
\end{remark}

% [Tutor Note: 可核对的要点有二。（i）$\tau^{ij}=\partial a_D^i/\partial a_j$（本页起首）与式 (20.16) 之后的 $\tau^{ij}=\partial^2\mathcal{F}/\partial a_i\partial a_j$ 一致，因为 $a_D^i=\partial\mathcal{F}/\partial a_i$。（ii）Riemann 双线性关系给出 $(v_{ij}/u_{ij})$ 虚部正定，于是式 (20.16) 的 $\dd s^2=\operatorname{Im}\tau^{ij}\dd a_i\dd\bar a_j$ 是正定厄米度量（$\tau^{ij}$ 对称 $\Rightarrow$ $\operatorname{Im}\tau$ 对称），这正是「$\tau$ 非退化且虚部正定」这一约束的几何含义。]
\subsection{可积系统}

\begin{definition}[20.34]
规范群为 $G$ 的 $\mathcal{N}=2^*$ 理论是由一个 $G$-矢量超场与一个位于 $G$ 的伴随表示中的有质量超多重态构成的拉格朗日 $4d$ $\mathcal{N}=2$ 理论。把该超多重态的质量取为零，即得到规范群为 $G$ 的 $4d$ $\mathcal{N}=4$ SYM 理论。
\end{definition}

\begin{remark}[20.35]
与 $\mathcal{N}=4$ SYM 理论一样，视规范群 $G$ 而定，$\mathcal{N}=2^*$ 理论可能存在不同的整体变体。
\end{remark}

$\mathcal{N}=2^*$ 理论也可以在 Seiberg–Witten 框架内描述，不过此时的 Seiberg–Witten 曲线不再是超椭圆曲线（纯 SYM 与带基本超多重态的 SYM 是超椭圆的），而是带穿孔环面的分歧覆叠 [114]。Donagi 与 Witten 对 $\mathcal{N}=2^*$ 理论的研究，使人们对描述 $4d$ $\mathcal{N}=2$ 量子场论低能物理的底层数学结构获得了深刻洞见。

\begin{proposition}[{{{{[114]}}}}]
一般而言，任意秩 $r$ 的 $4d$ $\mathcal{N}=2$ 理论在库仑分支上的低能有效拉格朗日量定义一个紧复 $r$ 维环面 $X_{\mathbf{u}}$ 的族 $X\to\mathcal{C}$，即阿贝尔簇。如

% [原文 p.316]
同前面的情形，$\mathcal{C}$ 上的特殊坐标连同质量参数，可以通过取 $X$ 上亚纯微分 $\lambda$ 的周期积分得到，其中 $\lambda$ 限制到纤维化 $X\to\mathcal{C}$ 的纤维上时是闭的，且定义到相差恰当 1-形式。$\lambda$ 的外微分 $\dd\lambda$ 是 $X$ 上无歧义的非退化、闭、全纯 2-形式 $\omega$，它限制到纤维上为零。
\end{proposition}

DSZ 配对自然在纤维 $X_{\mathbf{u}}$ 上诱导一个极化，从而给出 $\mathcal{C}$ 上一族带全纯 2-形式 $\omega$ 的极化阿贝尔簇。该形式定义一个复辛结构，从而是 $X$ 上的一个泊松结构。库仑分支坐标 $u_1,\dots,u_r$ 关于这个泊松括号构成一组极大的、相互交换且代数独立的函数。换言之：

\begin{proposition}[20.37]
任意秩 $r$ 的 $4d$ $\mathcal{N}=2$ 理论在库仑分支上的低能有效拉格朗日量定义一个复完全可积系统。
\end{proposition}

刻画 $4d$ $\mathcal{N}=2$ 理论低能物理的可积系统的交换哈密顿量，是低能 $\mathcal{N}=2$ 阿贝尔矢量场中标量场的 VEV。它们生成的轨道编码低能阿贝尔矢量超场的动能。

\begin{example}[20.38]
在前面讨论的情形中，Seiberg–Witten 曲线是黎曼球面 $\C\PP^1$ 的分歧覆叠，相应的阿贝尔簇是 Seiberg–Witten 曲线 $\Sigma_{\mathbf{u}}$ 的雅可比簇 $J(\Sigma_{\mathbf{u}})$。
\end{example}

存在一大类复完全可积系统，称为 Hitchin 可积系统 [198, 197]，它涵盖前面讨论的全部情形以及更一般的情形，例如 $4d$ $\mathcal{N}=2^*$ 理论。Hitchin 系统与亏格为 $g$ 的任意黎曼曲面 $C$ 以及约化紧李群 $G$ 相伴，它们来自在 $C$ 上规范群为 $G$ 的 $4d$ 自对偶（瞬子）Yang–Mills 方程的维数约化。用更偏高能物理的语言，Hitchin 系统由 $2d$ 伴随 QCD 理论构造。非紧黎曼曲面上的 Hitchin 系统在 [259, 113] 中构造。关于 Hitchin 可积系统的一般方面，读者可参见 Hitchin 的原始论文 [198, 197] 与专著 [196]。

Hitchin 系统中的一个点由 $C$ 上的全纯 $G_{\C}$-丛 $P_{G_{\C}}$（其中 $G_{\C}$ 是实李群 $G$ 的复化）以及 $C$ 上的伴随值全纯 1-形式 $\phi$ 组成。这一系统的可积结构在引入谱曲线或谱覆盖的概念后变得更加显式。可积系统结构涉及谱曲线 $\Sigma$ 的雅可比簇。

\begin{definition}[20.39]
$C$ 的谱覆盖是由对 $(P_{G_{\C}},\phi)$ 决定的 $C$ 的分歧覆叠 $\Sigma\subset T^*C$，其中 $T^*C$ 是 $C$ 的余切丛。
\end{definition}

对 $G=\SU(N)$，其中 $G_{\C}=\SL_N(\C)$，谱覆盖 $\Sigma$ 由行列式方程定义：

% [原文 p.317]
\begin{equation}
\det(\eta-\phi)=\eta^{\otimes N}+\eta^{\otimes(N-2)}W_2(\phi)+\eta^{\otimes(N-3)}W_3(\phi)+W_N(\phi)=0,\tag{20.19}
\end{equation}
其中 $\eta$ 是余切丛 $T^*C$ 上的重言 1-形式。这个方程一般定义 $T^*C$ 中一个亏格为 $\tilde{g}=(N^2-1)(g-1)+g$ 的黎曼曲面 $\Sigma$，它是 $C$ 的 $N$ 次分歧覆叠。$C$ 上的全纯 $i$-微分 $W_i(\phi)$ 由 $\phi$ 决定。

% [Tutor Note: 式 (20.19) 中缺少 $\eta^{\otimes(N-1)}W_1(\phi)$ 项，这并非笔误：对 $G=\SU(N)$，伴随表示中的 $\phi$ 无迹，故 $W_1(\phi)=\tr\phi=0$，该项消失。一般地 $\det(\eta-\phi)=\sum_{k=0}^{N}\eta^{\otimes(N-k)}(-1)^kW_k(\phi)$（各文献对 $W_k$ 的符号与归一化约定不同），式 (20.19) 只显式写出 $k=0,2,3,N$ 四项；$N=3$ 时中间两项合并为一项，故该式应按「逐项列出相异项」来理解。核对：取 $N=4$，$\phi$ 的特征值为 $\lambda_1,\dots,\lambda_4$，由 $\sum_i\lambda_i=0$ 得 $\det(\eta-\phi)=\prod_i(\eta-\lambda_i)=\eta^4+e_2\eta^2-e_3\eta+e_4$，与式 (20.19) 的 $\eta^{\otimes4}+\eta^{\otimes2}W_2+\eta W_3+W_4$ 逐项对应，即 $W_2=e_2$、$W_3=-e_3$、$W_4=e_4$。]
Hitchin 系统结构 $X\to\mathcal{B}$（其中纤维是阿贝尔簇）如下：底 $\mathcal{B}$ 是微分 $W_i(\phi)$ 构成的空间，形成维数为 $r=(N-1)$ 的复仿射空间。每条纤维是与 $\mathcal{B}$ 上一点相伴的谱曲线 $\Sigma$ 的雅可比簇。还存在一个特殊的截面：$C$ 上的一般点对应于 $\Sigma$ 上的 $N$ 个不同点，即式 (20.19) 的特征值。这些特征值定义 $\Sigma$ 上的一个全纯线丛，从而是雅可比簇 $J(\Sigma)$ 中的一个点。

\begin{example}[20.40]
对 SYM 理论（纯理论或带基本物质），曲线 $C$ 是黎曼球面 $\C\PP^1$。对 $\mathcal{N}=2^*$ 理论，$C$ 是带穿孔的环面。群 $G$ 是所考虑的 $4d$ $\mathcal{N}=2$ 理论的规范群。Hitchin 系统的底 $\mathcal{B}$ 与库仑分支 $\mathcal{C}$ 等同，全纯 $k$-微分 $W_k(\phi)$ 对应于库仑分支算子。Seiberg–Witten 曲线 $\Sigma$ 与由 $\mathcal{C}$ 上一点决定的谱覆盖等同。
\end{example}

刻画纯 SYM 理论或 $\mathcal{N}=2^*$ 理论低能物理的 Hitchin 系统，对应于众所周知的复完全可积系统族。纯 $\SU(N)$ SYM 理论与所谓的周期 Toda 链相伴，而 $\mathcal{N}=2^*$ 理论对应于椭圆 Calogero–Moser 系统。关于这些可积系统与 Seiberg–Witten 理论之间联系的更多细节，我们参见 [106] 中的教学式导引及其中的参考文献。

\begin{remark}[20.41]
这一讨论引出如下问题：是否所有 Hitchin 系统都刻画了特定 $4d$ $\mathcal{N}=2$ 理论在库仑分支上的低能物理？我们将在下一章深入探讨并给出答案。
\end{remark}

\subsection{值得注意的 $4d$ $\mathcal{N}=2$ SCFT}

让我们强调一些有趣的、略带异域色彩的理论，它们是通过 Seiberg 与 Witten 所发展的形式体系被发现的。

一类值得注意的例子是所谓的 Argyres–Douglas 理论 [24, 27]。这些 $4d$ $\mathcal{N}=2$ 超共形场论（SCFT）作为 $4d$ $\mathcal{N}=2$ 理论库仑分支 $\mathcal{C}$ 上特定奇点处的低能描述而出现，特别是当覆叠 $\Sigma\to\mathcal{C}$ 的多个分支点相撞时。一个原型情形出现在带一个有质量基本超多重态的 $SU(2)$ SYM 中。这一设定的 Seiberg–Witten 曲线由式 (20.13) 给出，我们在此重录：

% [原文 p.318]
\[
(\Sigma_{u,m}):\quad \Lambda^2 z+\frac{2\Lambda(x-m)}{z}=x^2-u,
\]
其中 $u\in\mathcal{C}\simeq\C$，$\Lambda$ 是动力学标度，$m$ 是基本超多重态的裸质量，$z\in\C\PP^1$ 且 $x\in\C$。这条 Seiberg–Witten 曲线 $\Sigma_{u,m}$ 是 $\C\PP^1_z$ 的 2 次分歧覆叠，其分歧点由判别式方程
\[
\Delta=z^3+\frac{uz^2}{\Lambda^2}-\frac{2mz}{\Lambda}+1=0,
\]
连同点 $z=\infty$ 确定。$\mathcal{C}$ 中的奇点在两个分支点相撞时出现，具体地，出现在使判别式的判别式
\[
\Delta'=\Delta(\Delta)=u^3-m^2u^2+9\Lambda^3mu+\frac{27}{4}\Lambda^6-8\Lambda^3m^3
\]
作为 $u$ 的函数取零的那些 $u$ 值处。如所预期，在 $\C$ 中有三个解。$\Delta'$ 的判别式为
\[
\Delta(\Delta')=\Delta(\Delta(\Delta))=8m^3+27\Lambda^3.
\]
于是，当 $m^3=-27\Lambda^3/8$ 时，两个奇点相撞，导致一个更复杂的奇点，在那里 Seiberg–Witten 曲线的两个独立圈收缩，产生尖点奇点（与简单奇点即节点相对照）。

值得注意的是，在 $\mathcal{C}$ 中这样的特定奇点处，低能理论包含相互非局域的无质量粒子，例如夸克与磁单极子，这表明该 SCFT 不能用拉格朗日量描述。关于 Argyres–Douglas 理论及其性质的进一步探讨，我们请读者参见 [339, 第 10 章] 及其中的参考文献。

另一组重要的 $4d$ $\mathcal{N}=2$ SCFT 是通过 Seiberg–Witten 理论预言其存在的所谓 Minahan–Nemeschansky 例外理论 [266, 267]。它们在这类理论的图景中扮演着显著角色。

在 Seiberg–Witten 理论的语境中，秩 1 的 $4d$ $\mathcal{N}=2$ SCFT 作为秩 1 的 $4d$ $\mathcal{N}=2$ 系统库仑分支 $\mathcal{C}$ 上奇点附近物理的低能描述而出现 [27]。远离奇点时，低能行为由 $4d$ $\mathcal{N}=2$ Maxwell 理论支配。$\mathcal{C}$ 上最简单的奇点类型是节点型的，即 Seiberg–Witten 曲线在那里出现二重点。在这样的奇点处，低能理论是 $4d$ $\mathcal{N}=2$ Maxwell 理论与无质量带电超多重态耦合。

然而，也可能出现更复杂的奇点，例如多个更简单奇点合并时。我们已经看到，在这些更复杂的奇点附近，低能动力学由 Argyres–Douglas 理论描述。通向秩 1 $4d$ $\mathcal{N}=1$ SCFT 的另一条途径是考虑无质量 $4d$ $\mathcal{N}=2$ $SU(2)$ 超对称 Yang–Mills 理论（带 $F\ge4$ 个基本味）库仑分支原点处的低能理论。具体而言，当 $F=4$ 时理论是共形的，而当 $F>4$ 时理论变为红外自由的。

% [Tutor Note: 两处可核对/存疑。（i）原文「通向秩 1 $4d$ $\mathcal{N}=1$ SCFT 的另一条途径」中的 $\mathcal{N}=1$ 按上下文应为 $\mathcal{N}=2$：该句讨论无质量 $4d$ $\mathcal{N}=2$ $SU(2)$ SYM 带 $F\ge4$ 个基本味时库仑分支原点处的低能理论，而 $F=4$ 共形、$F>4$ 红外自由正是 $\mathcal{N}=2$ 的行为（本节标题本身即「Remarkable $4d$ $\mathcal{N}=2$ SCFTs」）。此处照原文译出。（ii）判别式链可核：$\Delta$ 关于 $z$ 是三次的，$\Delta'=\Delta(\Delta)$ 关于 $u$ 是三次的，故「如所预期，在 $\C$ 中有三个解」自洽；再由 $\Delta(\Delta')=8m^3+27\Lambda^3=0$ 得 $m^3=-27\Lambda^3/8$，对应两个奇点合并（此时 $\Delta'$ 有重根）。]
% [原文 p.319]
描述奇点附近物理的 SCFT 只依赖于奇点附近的局部几何特征，这导致对可能出现的奇点进行分类的问题。在 Seiberg–Witten 理论中，每个秩 1 的 $4d$ $\mathcal{N}=2$ 理论都与库仑分支上复环面的纤维化 $\Sigma\to\mathcal{C}$ 相伴。该纤维化的全空间构成一个复二维代数簇。这类簇的所有可能奇点已由 Kodaira 分类 [236]，它们在局部由商 $\C^2/\Gamma$ 建模，其中 $\Gamma$ 是 $\SU(2)$ 的有限子群。

$\SU(2)$ 的有限子群 $\Gamma$ 具有 $ADE$ 分类。值得注意的是，对应于 $\C^2/\Gamma$ 的 $4d$ $\mathcal{N}=2$ SCFT 的味对称李群的 Dynkin 型与 $\Gamma$ 的型一致。虽然 $A$ 型与 $D$ 型对应于前面讨论的理论——例如带物质的 $4d$ $\mathcal{N}=2$ Maxwell 理论、Argyres–Douglas 理论，或带 $F\ge4$ 个基本味的 $SU(2)$ SYM——但例外情形 $E_6$、$E_7$ 与 $E_8$ 引出了新的 $4d$ $\mathcal{N}=2$ SCFT 的发现。这些 SCFT 分别具有 $E_6$、$E_7$ 与 $E_8$ 型的味对称群，被称为 Minahan–Nemeschansky 理论 [266, 267]。

\subsection{瞬子计数}

我们以强调一个与 Seiberg–Witten 理论同时代的重大进展来结束本章，即瞬子计数纲领。Seiberg–Witten 理论的一个核心成就是间接确定了低能预势，而低能预势又编码了规范耦合重整化的非微扰瞬子修正。\footnote{凭借 $\mathcal{N}=2$ 超对称，全纯规范耦合与物理规范耦合一致，而全纯性把复化规范耦合的重整化限制为只由单圈微扰贡献与非微扰瞬子效应组成（18.3.1 节）。} 关于纯 $\mathcal{N}=2$ $SU(2)$ SYM 情形的显式讨论，参见 [245]。

自然要问：由 Seiberg–Witten 理论确定的预势，能否通过对瞬子构型上的路径积分更直接地导出？这个问题推动了多瞬子模空间（即瞬子数为 $n>1$ 的构型）研究的重大进展，最终汇成综述 [118]。现代瞬子计数纲领由 Nekrasov 的开创性工作 [284] 发起，其中证明 [118] 的多瞬子积分（它们给出 Seiberg–Witten 理论的低能预势）可以解释为等变积分，从而用等变局部化技术计算。

在物理上，直接给出 $\mathcal{N}=2$ 规范理论低能预势的等变积分，可用理论的一个特定形变来自然解释，该形变由两个实常数 $\epsilon_1$ 与 $\epsilon_2$ 参数化。具体而言，考虑与给定 $4d$ $\mathcal{N}=2$ 规范理论相伴的 $5d$ $\mathcal{N}=1$ 理论，把它放在半径为 $\beta$ 的圆周 $S^1_\xi$ 上的 $\C^2_{z_1,z_2}$ 纤维化上，并作如下等同：
\[
(z_1,z_2,\xi)\sim\big(\ee^{\ii\beta\epsilon_1}z_1,\ee^{\ii\beta\epsilon_2}z_2,\xi+\beta\big).
\]

% [原文 p.320]
所得几何称为 $\Omega$-背景。在这一设定下，超对称配分函数 $Z^{5d}(\beta,\epsilon_1,\epsilon_2,\dots)$ 具有良定义的四维极限
\[
Z^{4d}(\epsilon_1,\epsilon_2,\dots)=\lim_{\beta\to0}Z^{5d}(\beta,\epsilon_1,\epsilon_2,\dots),
\]
其中省略号表示当规范理论存在整体对称性时与这些对称性相伴的可能 fugacity。关键在于，这个配分函数编码非微扰低能动力学：Seiberg–Witten 预势在如下极限中恢复
\[
\mathcal{F}(\dots)=\lim_{\epsilon_1,\epsilon_2\to0}\epsilon_1\epsilon_2\log Z^{4d}(\epsilon_1,\epsilon_2,\dots).
\]
超对称配分函数 $Z^{4d}(\epsilon_1,\epsilon_2,\dots)$ 可对一大类规范理论计算，包括具有任意单规范群与物质内容的 $\mathcal{N}=2$ SYM 理论，以及酉箭图规范理论。

我们将在本书中不再深入瞬子计数这一主题，而请读者参见综述 [326, 54, 340]，它们详细考察了这一主题并强调其与规范理论和弦理论中广泛进展的联系。一个尤其值得注意的进展是 AGT 对应（综述见 [244]），它在四维 $\mathcal{N}=2$ 的 S 类理论（将在第 21 章讨论）与二维共形场论之间建立了深刻联系。W-代数（它们本身与高阶泰希米勒理论 [55] 有关）在这一语境中的出现，在 [340] 中有综述。与本书主题更直接相关的是瞬子计数与谱网络通过精确 WKB 方法的相互交织，这已在近期工作 [208] 中探讨。

% [Tutor Note: $\Omega$-背景的等同关系 $(z_1,z_2,\xi)\sim(\ee^{\ii\beta\epsilon_1}z_1,\ee^{\ii\beta\epsilon_2}z_2,\xi+\beta)$ 的含义是：$5d$ 理论放在 $\C^2_{z_1,z_2}$ 的纤维化上、底为半径 $\beta$ 的 $S^1_\xi$，绕 $S^1$ 一周时两个复平面分别整体转动 $\beta\epsilon_1$、$\beta\epsilon_2$。$\beta\to0$ 时 $S^1$ 收缩，得到四维配分函数 $Z^{4d}$；而 $\epsilon_1\epsilon_2\log Z^{4d}$ 的 $\epsilon_i\to0$ 极限给出预势 $\mathcal{F}$，因为规范理论配分函数按 $Z^{4d}\sim\exp\big(\mathcal{F}/(\epsilon_1\epsilon_2)+\cdots\big)$ 展开。核对：把该展开代入极限式，非微扰部分（瞬子项）的系数恰为 $\mathcal{F}$ 的瞬子展开系数；符号约定若取 $\epsilon_i\to-\epsilon_i$，$\mathcal{F}$ 可能差一个整体符号，视文献而定。]
